\documentclass[12pt,a4paper]{ctexart}

\usepackage{geometry}
\usepackage{booktabs}
\usepackage{tabularx}
\usepackage{array}
\usepackage{xcolor}
\usepackage{longtable}
\usepackage{amsmath}

\geometry{
    left=2.2cm,
    right=2.2cm,
    top=2.2cm,
    bottom=2.2cm
}

\setlength{\parindent}{0pt}
\renewcommand{\arraystretch}{1.35}

\definecolor{headblue}{RGB}{42,92,140}
\definecolor{lightblue}{RGB}{235,243,250}

\newcolumntype{Y}{>{\centering\arraybackslash}X}
\newcolumntype{L}{>{\raggedright\arraybackslash}X}

\begin{document}

\begin{center}
    {\LARGE\bfseries STM32F411RE 外设接口分配表}\\[0.5em]
    {\large 基于 NUCLEO-F411RE 的空气质量监测系统}
\end{center}

\vspace{1em}

\section{HC-06 蓝牙模块}

HC-06 通过 STM32F411RE 的 USART2 与主控通信。

\begin{table}[htbp]
\centering
\rowcolors{2}{lightblue}{white}
\begin{tabularx}{0.9\textwidth}{>{\bfseries}Y Y Y}
\toprule
\rowcolor{headblue}
\color{white}HC-06 信号 &
\color{white}STM32F411RE 引脚 &
\color{white}STM32 功能 \\
\midrule
HC06\_TXD & PA3 & USART2\_RX \\
HC06\_RXD & PA2 & USART2\_TX \\
VCC & NUCLEO\_5V & 5 V 供电 \\
GND & GND & 系统地 \\
\bottomrule
\end{tabularx}
\end{table}

串口采用交叉连接：

\[
\text{HC-06 TXD}
\rightarrow
\text{STM32 PA3 / USART2\_RX}
\]

\[
\text{HC-06 RXD}
\leftarrow
\text{STM32 PA2 / USART2\_TX}
\]

因此 CubeMX 中配置为：

\begin{verbatim}
PA2 = USART2_TX
PA3 = USART2_RX
\end{verbatim}

\section{2.8 英寸 ILI9341 TFT LCD}

2.8 英寸 TFT LCD 使用 SPI1 与 STM32F411RE 通信，
显示控制器为 ILI9341，分辨率为 $240\times320$。

\begin{table}[htbp]
\centering
\rowcolors{2}{lightblue}{white}
\begin{tabularx}{0.95\textwidth}{>{\bfseries}Y Y Y}
\toprule
\rowcolor{headblue}
\color{white}TFT 信号 &
\color{white}STM32F411RE 引脚 &
\color{white}功能 \\
\midrule
TFT\_SCK & PA5 & SPI1\_SCK \\
TFT\_SDO & PA6 & SPI1\_MISO \\
TFT\_SDI & PA7 & SPI1\_MOSI \\
TFT\_CS & PB6 & GPIO Output \\
TFT\_RS / DC & PC7 & GPIO Output \\
TFT\_RESET & PB4 & GPIO Output \\
TFT\_BL & PB3 & GPIO Output \\
VCC & NUCLEO\_5V & 5 V 供电 \\
GND & GND & 系统地 \\
\bottomrule
\end{tabularx}
\end{table}

SPI1 的核心连接关系为：

\begin{verbatim}
PA5 -> TFT_SCK
PA6 -> TFT_SDO / MISO
PA7 -> TFT_SDI / MOSI

PB6 -> TFT_CS
PC7 -> TFT_RS / DC
PB4 -> TFT_RESET
PB3 -> TFT_BL
\end{verbatim}

本项目当前仅实现显示功能，因此触摸屏相关引脚暂不使用：

\begin{verbatim}
T_CLK
T_CS
T_DIN
T_DO
T_IRQ
\end{verbatim}

\section{PMS5003ST 多参数环境传感器}

PMS5003ST 通过 USART1 与 STM32F411RE 通信，
同时使用两个 GPIO 控制 SET 和 RESET 信号。

\begin{table}[htbp]
\centering
\rowcolors{2}{lightblue}{white}
\begin{tabularx}{0.95\textwidth}{>{\bfseries}Y Y Y}
\toprule
\rowcolor{headblue}
\color{white}PMS5003ST 信号 &
\color{white}STM32F411RE 引脚 &
\color{white}STM32 功能 \\
\midrule
PMS\_TXD & PA10 & USART1\_RX \\
PMS\_RXD & PA9 & USART1\_TX \\
PMS\_SET & PB10 & GPIO Output \\
PMS\_RESET & PB5 & GPIO Output \\
VCC & NUCLEO\_5V & 5 V 供电 \\
GND & GND & 系统地 \\
\bottomrule
\end{tabularx}
\end{table}

串口同样采用交叉连接：

\[
\text{PMS TXD}
\rightarrow
\text{STM32 PA10 / USART1\_RX}
\]

\[
\text{PMS RXD}
\leftarrow
\text{STM32 PA9 / USART1\_TX}
\]

CubeMX 中配置为：

\begin{verbatim}
PA9  = USART1_TX
PA10 = USART1_RX

PB10 = GPIO_Output  // PMS_SET
PB5  = GPIO_Output  // PMS_RESET
\end{verbatim}

\section{系统引脚汇总}

\begin{table}[htbp]
\centering
\rowcolors{2}{lightblue}{white}
\begin{tabularx}{\textwidth}{>{\bfseries}Y L}
\toprule
\rowcolor{headblue}
\color{white}模块 &
\color{white}STM32F411RE 引脚分配 \\
\midrule
HC-06 &
PA2 -- USART2\_TX；
PA3 -- USART2\_RX \\

PMS5003ST &
PA9 -- USART1\_TX；
PA10 -- USART1\_RX；
PB10 -- PMS\_SET；
PB5 -- PMS\_RESET \\

ILI9341 SPI &
PA5 -- SPI1\_SCK；
PA6 -- SPI1\_MISO；
PA7 -- SPI1\_MOSI \\

ILI9341 控制 &
PB6 -- TFT\_CS；
PC7 -- TFT\_RS/DC；
PB4 -- TFT\_RESET；
PB3 -- TFT\_BL \\
\bottomrule
\end{tabularx}
\end{table}

\section{STM32F411RE 核心接口一览}

最终系统主要使用以下 13 个 MCU 引脚：

\begin{center}
\begin{tabular}{cccc}
\toprule
\textbf{引脚} & \textbf{功能} &
\textbf{引脚} & \textbf{功能} \\
\midrule
PA2  & HC-06 TX &
PA3  & HC-06 RX \\

PA5  & TFT SCK &
PA6  & TFT MISO \\

PA7  & TFT MOSI &
PA9  & PMS TX \\

PA10 & PMS RX &
PB3  & TFT BL \\

PB4  & TFT RESET &
PB5  & PMS RESET \\

PB6  & TFT CS &
PB10 & PMS SET \\

PC7  & TFT DC &
      &          \\
\bottomrule
\end{tabular}
\end{center}

\vfill

\begin{center}
\small
NUCLEO-F411RE / STM32F411RE \quad
PMS5003ST + ILI9341 + HC-06
\end{center}

\end{document}